\chapter*{Glossary}
\addcontentsline{toc}{chapter}{Glossary}

\begin{description}[leftmargin=!,labelwidth=4.5cm,itemsep=0.5em]
\item[Abdominal thrusts (Heimlich)] Inward-upward thrusts into the upper abdomen used to expel an airway obstruction in conscious adults and children over 1 year.
\item[AED] Automated External Defibrillator -- a portable device that analyzes heart rhythm and delivers an electric shock if needed.
\item[Agonal breathing] Occasional, gasping, irregular breaths in a person in cardiac arrest. This is NOT normal breathing -- begin CPR.
\item[Airway] The passage through which air reaches the lungs (mouth, nose, throat, windpipe).
\item[Anaphylaxis] A rapid, life-threatening whole-body allergic reaction requiring epinephrine and emergency care.
\item[Arterial bleeding] Bright-red blood that spurts in time with the heartbeat; can be rapidly fatal.
\item[AVPU] Alert, Voice, Pain, Unresponsive -- a scale for rating level of consciousness.
\item[Avulsion] A wound in which skin or tissue is partially or completely torn away.
\item[Brachial artery] The pulse point on the inner upper arm; used for infants and children.
\item[Capillary refill] The time for color to return to a pressed nail bed (normally under 2 seconds).
\item[Cardiac arrest] Sudden loss of heart function and breathing; treated with CPR and defibrillation.
\item[Carotid artery] The main pulse point in the neck, beside the windpipe.
\item[CPR] Cardiopulmonary resuscitation -- chest compressions combined with rescue breaths.
\item[Contusion] A bruise; damaged blood vessels beneath intact skin.
\item[Cyanosis] Bluish/gray discoloration of skin or lips from low blood oxygen.
\item[DCAP-BTLS] Deformities, Contusions, Abrasions, Punctures, Burns, Tenderness, Lacerations, Swelling -- what to look for in a head-to-toe survey.
\item[Defibrillation] Delivery of an electric shock to restore a normal heart rhythm.
\item[DKA] Diabetic ketoacidosis -- a life-threatening complication of high blood sugar (type 1 diabetes).
\item[Epistaxis] Nosebleed.
\item[Epinephrine] A medication (adrenaline) used to reverse anaphylaxis; given by auto-injector.
\item[FAST] Face, Arm, Speech, Time -- the stroke recognition mnemonic.
\item[Febrile] Relating to fever.
\item[Frostbite] Freezing of skin and deeper tissues from cold exposure.
\item[Glasgow Coma Scale (GCS)] A standardized assessment (Eye, Verbal, Motor responses) scored 3-15 to quantify level of consciousness after head injury.
\item[Good Samaritan law] Legal protection for those who voluntarily give reasonable first aid in good faith.
\item[Hemorrhage] Bleeding; loss of blood from the circulation.
\item[Hypoglycemia] Low blood sugar; a diabetic emergency.
\item[Hyperglycemia] High blood sugar.
\item[Hypothermia] Dangerously low body temperature (below 95\textdegree F / 35\textdegree C).
\item[Hypovolemic shock] Shock caused by loss of blood or fluid volume.
\item[Ischemia] Reduced blood flow to a tissue, causing oxygen deprivation.
\item[Laceration] A cut or tear in the skin and underlying tissue.
\item[OPQRST] Onset, Provocation, Quality, Radiation, Severity, Time -- a system for describing pain or symptoms.
\item[Primary survey] The first rapid check for life threats (scene safety, responsiveness, airway/breathing, circulation).
\item[Pulse point] A place where an artery can be felt against bone (carotid, brachial, radial, femoral).
\item[Recovery position] A stable side-lying position that keeps the airway open in an unconscious, breathing person.
\item[RICE] Rest, Ice, Compression, Elevation -- for sprains, strains, and bruises.
\item[SAMPLE] Symptoms, Allergies, Medications, Past history, Last meal, Events -- a history-taking mnemonic.
\item[Secondary survey] A systematic head-to-toe examination to find all injuries after life threats are managed.
\item[Shock] A state in which vital organs do not receive enough oxygen-rich blood.
\item[Sprain] Stretching or tearing of a ligament.
\item[START] Simple Triage And Rapid Treatment -- a method for sorting casualties in a mass-casualty incident.
\item[Strain] Stretching or tearing of a muscle or tendon.
\item[Syncope] Fainting; a temporary loss of consciousness from reduced blood flow to the brain.
\item[Tachycardia] Abnormally fast heart rate.
\item[TIA] Transient ischemic attack (``mini-stroke''); temporary stroke-like symptoms that still require urgent medical evaluation.
\item[Triage] Sorting patients by urgency when resources are limited.
\item[Tourniquet] A tight band applied above severe extremity bleeding as a last resort.
\item[Venous bleeding] Darker blood that flows steadily from damaged veins.
\item[Conjunctival pallor] Paleness of the inner eyelid mucosa indicating anemia or poor perfusion.
\item[Crepitus] A grating sensation or sound caused by fractured bone ends rubbing together.
\item[Decerebrate posturing] Abnormal extension of all four limbs indicating severe brain injury.
\item[Decorticate posturing] Abnormal flexion of arms with extension of legs indicating brain injury.
\item[Evisceration] Protrusion of internal organs through an open wound.
\item[Heterochromia] Difference in coloration of the irises; may be normal or indicate injury.
\item[Jugular venous distension (JVD)] Bulging neck veins indicating increased central venous pressure or cardiac tamponade.
\item[Kussmaul breathing] Deep, rapid breathing associated with diabetic ketoacidosis.
\item[Mandibular fracture] Break of the jaw bone; suspect with facial trauma, malocclusion, or crepitus.
\item[Meniscus tear] Injury to the cartilage cushion in the knee joint; causes pain, swelling, and locking.
\item[Nasal septal hematoma] Blood collection between nasal cartilage and mucosa; requires drainage to prevent tissue death.
\item[Oliguria] Decreased urine output (less than 0.5 mL/kg/hr); indicates poor perfusion or shock.
\item[Pericardial tamponade] Fluid accumulation around the heart compressing it; causes JVD, muffled heart sounds, and hypotension.
\item[Phalangeal fracture] Break of a finger or toe bone; immobilize with buddy taping.
\item[Retrobulbar hematoma] Bleeding behind the eye causing bulging, vision loss, and increased intraocular pressure.
\item[Subconjunctival hemorrhage] Broken blood vessel on the white of the eye; appears dramatic but is usually benign.
\item[Tympanic membrane rupture] Torn eardrum from blunt trauma or pressure change; protect from water and refer to physician.
\item[Volar splint] Splint applied to the palm side of the wrist and hand for fractures or soft tissue injuries.
\item[In-line stabilization] Holding the head and neck in neutral alignment to protect the cervical spine.
\item[Anterior-posterior position] Alternative AED pad placement with one pad on the front chest and one on the back.
\end{description}

\index{glossary}
